\section{Fragestellung: ein zweiter, unabhängiger Zugang zu $\theta=2/9$}
Der Koide-Winkel $\theta=2/9$ ist empirisch auf sieben Stellen bestätigt und in FFGFT als
rationale Adresse eines über einen Referenzpunkt fixierten Wertes verstanden (Dok.~292), nicht
als Polmassen-Exaktheitsanspruch (P42). Dok.~290/291 haben gezeigt, in welchem \emph{Kanal}
$\theta$ lebt: nicht im symmetrischen Betragsfunktional (dort ist $2/9$ counter-blind), sondern
im Phasenkanal, mit $\chi=\pi/2$ als schief-adjungiertem Selektor und $\delta$ als noch offener
zweiter Harmonischer. Der \emph{Betrag} von $\theta$ blieb dabei eine gefundene, nicht
hergeleitete Zahl.

Dieses Dokument liefert einen zweiten, davon völlig unabhängigen Zugang. Er stammt nicht aus den
Massen und nicht aus der Koide-Parametrisierung, sondern rein aus der \emph{Geometrie} der
Verdopplung $6=3\oplus3'$ (Dok.~285): der Einbettung der FFGFT-eigenen $C_3$-Struktur in die
ikosaedrische Ambientgruppe $A_5$. Diese Einbettung erzeugt $\theta=2/9$ parameterfrei,
vorwärts und exakt -- ohne $2/9$ als Eingabe. Dass ein zweiter, methodisch fremder Weg auf
\emph{denselben} rationalen Wert trifft, ist die eigentliche Aussage: nicht eine Übersetzung des
einen in den anderen, sondern ein Konvergenzbeleg im Sinne der bewährten FFGFT-Logik, in der
$\alpha$, die Leptonmassen und Koide als \emph{Verifikationen} verschiedener Zugänge auf
dieselben Konstanten zeigen. Die Rechnungen liegen in \texttt{ikosaeder\_theta\_delta.py}
(numpy-only, Seed 20780458).

\section{Die $C_3$-in-$A_5$-Einbettung und die volle Umverteilung}
Die drei Lepton-Moden sind die Fourier-Eigenmoden des $Z_3$-Zirkulanten auf der
$\mathbb{C}^3$-Faser (Dok.~282/288): die triviale Mode $v_0=\tfrac{1}{\sqrt3}(1,1,1)$ und die
zwei nicht-trivialen, galois-konjugierten Moden $v_{1,2}=\tfrac{1}{\sqrt3}(1,\omega^{\pm1},
\omega^{\pm2})$ mit $\omega=e^{2\pi i/3}$. Das Elektron ist eine nicht-triviale Mode; welche der
beiden, entscheidet allein das Vorzeichen des Winkels, nicht der hier berechnete Wert.

Die Verdopplung stellt $C_3$ in die ikosaedrische Gruppe $A_5$; die fünfzählige Erzeugende ist
eine Drehung $R_5$ um $72^\circ$ um eine ikosaedrische $5$-fach-Achse (Standardeinbettung:
Achse $(0,1,\varphi)$ mit $\varphi=\tfrac{1+\sqrt5}{2}$). Da die $3$-zählige und die $5$-zählige
Achse verschieden orientiert sind, mischt $R_5$ die $C_3$-Moden. Die \emph{volle} Umverteilung
des Elektron-Mode $v_e=v_1$ hat die Gewichte $p_j=\lvert\langle v_j\vert R_5\,v_e\rangle\rvert^2$:
\[
\underbrace{p_0=\frac{2}{9}}_{\text{trivial }v_0},\qquad
\underbrace{p_2=\frac{5-3\varphi}{9}}_{\text{2.\ Harm.\ }v_2},\qquad
\underbrace{p_1=\frac{2+3\varphi}{9}}_{\text{bleibt }v_e},\qquad p_0+p_1+p_2=1.
\]
Alle drei Gewichte tragen den Nenner $9=3^2$ -- dieselbe $Z_3$-Struktur, in der auch $2/9$ als
Orbit-Adresse sitzt. Numerisch bestätigt (40 Stellen, mpmath): $p_0=0{,}2222\ldots=2/9$ exakt,
Differenz $<10^{-30}$.

\section{Der entscheidende Befund: $\theta=2/9$ als trivialer Anteil}
Der Anteil, der unter der ikosaedrischen $5$-fach in die \emph{triviale} Mode $v_0$ übergeht, ist
\[
\boxed{\;p_0=\big\lvert\langle v_0\,\vert\,R_5\,v_e\rangle\big\rvert^2=\frac{2}{9}\;}
\]
exakt und parameterfrei. In den Wert geht kein Massenverhältnis, keine Koide-Parametrisierung
und kein $2/9$ ein; er folgt allein aus den $C_3$-Moden und der ikosaedrischen Drehung. Der
zugehörige \emph{totale} Leak ist $p_0+p_2=(7-3\varphi)/9=0{,}23843$, der zweite,
parameterfreie Kandidat für die Amplitude $\delta$ der zweiten Harmonischen; er stimmt mit dem
aus $\theta=2/9$ rückgerechneten $\delta^\ast=0{,}23887$ auf $0{,}18\,\%$ überein (unterhalb der
$\tau$-limitierten Toleranz von $0{,}135\,\%$ nur knapp, $\sim1{,}3\sigma$). Betrag \emph{und}
Winkel entstammen damit \emph{derselben} $5$-fach-Umverteilung: $\theta$ ist der triviale, $\delta$
der totale Anteil. Die frühere getrennte Behandlung als einzelne Harmonische (erste für $\theta$,
zweite für $\delta$) war die zu enge Sicht; die volle Umverteilung trägt beide.

\section{Härtetests: die Spezifität der ikosaedrischen Geometrie}
Der Wert $2/9$ ist keine generische Eigenschaft irgendeiner Drehung, sondern der ikosaedrischen:
\begin{itemize}
\item \textbf{Achsenspezifität.} $200$ zufällige Drehachsen bei $72^\circ$ ergeben den
$v_0$-Leak im Bereich $[0{,}036;\,0{,}452]$; \emph{keine einzige} trifft $2/9$. Nur die
ikosaedrischen $5$-fach-Achsen liefern es.
\item \textbf{Winkel- und Ordnungsspezifität.} Der Winkel-Scan zeigt $2/9$ genau bei $72^\circ$;
die $n$-fach-Reihe liefert $2/9$ ausschließlich bei $n=5$ (bei $n=3$ steht $2/5$, bei $144^\circ$
steht $4/9$). Die rationale Sauberkeit tritt an den ausgezeichneten Winkeln auf.
\item \textbf{Modenunabhängigkeit.} Beide nicht-trivialen Moden $v_1,v_2$ geben $p_0=2/9$; die
triviale $v_0$ gäbe $4/9$. Das Ergebnis hängt nicht davon ab, welche der beiden konjugierten
Moden als Elektron gelesen wird.
\item \textbf{Achsenwahl (diskret, P3).} Von den sechs $5$-fach-Achsen geben drei ($+\varphi$-Zweig,
$3$-$5$-Achswinkel $37{,}38^\circ$) den Wert $2/9$, drei ($-\varphi$-Zweig, $79{,}19^\circ$) den
Wert $4/9$. Die physikalische Einbettung $C_3<A_5$ wählt den nächstliegenden $5$-fach, den
$+\varphi$-Zweig -- also $2/9$. Es ist eine diskrete Zweifachwahl, kein freier Parameter.
\end{itemize}

\section{Zwei Zeugen, ein Wert}
Die $\theta$-Frage hat damit zwei voneinander unabhängige Quellen, die exakt zusammenfallen:
\begin{center}
\begin{tabular}{@{}p{0.31\linewidth}p{0.31\linewidth}p{0.30\linewidth}@{}}
\toprule
\textbf{Zugang} & \textbf{Herkunft} & \textbf{Wert}\\
\midrule
Koide (phänomenologisch) & aus den gemessenen Lepton-Massen; keine Geometrie & $\theta=2/9$\\
$C_3$-in-$A_5$ (geometrisch) & aus der $5$-fach-Umverteilung der Moden; keine Massen, kein $2/9$-Input & $p_0=2/9$\\
\bottomrule
\end{tabular}
\end{center}
Beide Wege wissen nichts voneinander -- der eine kennt nur Massen, der andere nur Geometrie --
und treffen dennoch \emph{exakt} dieselbe rationale Zahl. Das ist ein Konvergenzbeleg, kein
Übersetzungsdefizit: es ist nicht erforderlich, dass sich der geometrische Weg in die
Koide-Kosinus-Parametrisierung zurückübersetzt. Genau wie ein hergeleiteter Wert nicht mit
seiner Messvorschrift identisch sein muss, um sie zu bestätigen, muss die parameterfreie
geometrische Vorhersage $2/9$ nicht die Gestalt des phänomenologischen Fitparameters annehmen,
um mit ihm übereinzustimmen. Dass der geometrische Wert \emph{nicht} mit $\theta$ variiert,
ist dabei das Kennzeichen einer Vorhersage, nicht ein Mangel: die Geometrie besitzt keinen
freien Parameter und nennt einen festen Wert.

\section{$\theta=2/9$ als Übersetzungs-Invariant zweier Modelle}
Der \enquote{zweite Zeuge} ist genauer zu fassen, damit er nicht als ontologische Behauptung
missverstanden wird. Es geht \emph{nicht} darum, dass eine geometrische Wahrheit eine
phänomenologische schlägt, und auch nicht darum, dass der Torus \enquote{eigentlich} der
Hilbertraum sei oder umgekehrt. Beide -- die Punkt-\slash Differentialgeometrie des $T^4$ und die
Spektralgeometrie des Hilbertraums $L^2(T^4)\otimes\mathbb{C}^3$ -- sind \emph{mathematische
Modelle}, die nachweislich und verlustfrei ineinander übersetzbar sind (Dok.~230/282): die
$Z_3$-Fasern, die $C_3$-Fourier-Moden, die $A_5$-Einbettung sind in beiden Sprachen dieselbe
Struktur, nur verschieden dargestellt.

Der eigentliche Befund ist daher ein methodologischer, kein ontologischer: $\theta=2/9$ ist ein
\emph{Invariant der Übersetzung}. Er erscheint in der spektralen Darstellung (als
Umverteilungsgewicht $p_0$) ebenso wie in der phänomenologischen Koide-Adresse, weil er der
gemeinsamen Struktur gehört und nicht einer einzelnen Darstellung. Genau das macht ihn belastbar:
Die Konstante fällt aus der Struktur, unabhängig davon, in welcher der übersetzbaren Sprachen man
rechnet. Dies liegt auf der bewährten Linie, in der $\hbar$ und $c$ Umrechnungsfaktoren sind und
nicht ontologische Primitive: Was zählt, ist das unter der Umrechnung Invariante, nicht die Wahl
einer Seite als \enquote{wahr}. Über die Frage, welches der beiden Modelle die Wirklichkeit ist,
wird hier bewusst nichts behauptet -- sie wird nicht gebraucht, und die Rechnung liefert sie
nicht. Der Wert des Ergebnisses liegt gerade in seiner \emph{Modellunabhängigkeit}.

\section{Skaleninvarianz des Musters gegenüber der Besetzung}
In den geometrischen Wert $p_0=2/9$ und die Verhältnisse geht \emph{keine} Skala ein: Die
$C_3$-in-$A_5$-Geometrie kennt weder Masse noch Energie, sie liefert reine Zahlen -- den Winkel
und die Nenner-$9$-Gewichte. Damit ist die \emph{Musterschicht} -- Winkel plus Verhältnisse --
eine Eigenschaft der Topologie, nicht eines bestimmten Energiebereichs. Sie gilt überall dort,
wo die $Z_3$/$A_5$-Verdopplung auftritt; der Lepton-Sektor ist \emph{eine} Realisierung, nämlich
diejenige, in der die Skala $E_0$ (Dok.~292) die absolute Lage genau auf die gemessenen
$e,\mu,\tau$ legt.

Davon strikt zu trennen ist die \emph{Besetzungsschicht}: ob auf einer gegebenen Skala Zustände
tatsächlich sitzen. Die Geometrie sagt \enquote{\emph{falls} hier ein $Z_3$-Triplett realisiert
ist, dann mit diesem $\theta$ und diesen Verhältnissen} -- sie sagt nicht \enquote{hier ist
eines}. Auf einer anderen Skala steht dasselbe Zirkulant-Muster, ohne dass dort besetzte Zustände
liegen müssen. Das ist dieselbe Logik wie beim Turm $D(4+3n)$ und dem Leerheitssatz (Dok.~285):
das Muster ist überall präsent, die Besetzung ist eine separate, skalengebundene und
dynamikabhängige Frage; ein leeres Fenster trägt dasselbe Muster wie ein besetztes, nur ohne
Teilchen darin. Aus der Skaleninvarianz des Musters folgt daher \emph{keine} Vorhersage neuer
Teilchen auf anderen Skalen -- die Universalität betrifft die Struktur, nicht die Besetzung.

\section{Status und offene Kante}
Was gesichert ist: Die ikosaedrische Verdopplungsstruktur erzeugt $\theta=2/9$ parameterfrei,
vorwärts, exakt und ikosaeder-spezifisch, in unabhängiger Übereinstimmung mit dem
phänomenologischen Koide-Wert; die Amplitude $\delta=(7-3\varphi)/9$ entstammt derselben
Umverteilung. Damit hat die lange offene Frage -- ob die $T^4/Z_3$-Struktur $\theta=2/9$
erzwingt -- erstmals eine konkrete, exakte, geometrische Antwort im Rahmen der bereits
etablierten $A_5$-Verdopplung (Dok.~285).

Was als Ableitungsschritt offen bleibt (und ausdrücklich nicht als Einwand gegen den
Konvergenzbeleg): der explizite Mechanismus, der die geometrischen Umverteilungsgewichte
$(2/9,\,(5-3\varphi)/9,\,(2+3\varphi)/9)$ mit den physikalischen Massenverhältnissen verbindet.
Diese Brücke muss nicht über Koide laufen, sondern kann direkt über die Massen oder die
Modenstruktur geführt werden; sie ist der Gegenstand der Weiterarbeit. Die operative
Rolle der $5$-fach im Phasenkanal ist konsistent mit Dok.~291: $R_5$ setzt eine Phase $\pi/2$,
den Wert, den Dok.~291 mit $\chi$ verbindet.
